% Vietnamese translation for OI Public Library
% Source-File: IOI中国国家候选队论文/2006/陈启峰/附件/友好的动物的解题报告.doc
\documentclass[11pt]{article}
\usepackage{oipl-vn}

\title{Lời giải bài ``Những loài vật thân thiện''}
\author{Chen Qifeng}
\date{}

\begin{document}
\maketitle

\origtitle{Friendly Animals solution report}
\sourcepath{IOI 2006 / Chen Qifeng / attachments / Friendly Animals solution report.doc}

\section{Mô tả bài toán}

Hành tinh \(W\) có khí hậu thích hợp và hệ sinh vật phong phú như Trái Đất.
Sau nhiều năm nghiên cứu, các nhà sinh vật học ngoài hành tinh đã phát hiện
được hàng chục nghìn loài, và con số này vẫn tiếp tục tăng.

Sinh vật trên hành tinh \(W\) có quan hệ rất thú vị. Có những cặp loài rất
thân thiện, luôn ở gần nhau; nhưng cũng có những cặp loài thù địch, cứ gặp
nhau là dễ xảy ra tranh đấu. Để hiểu rõ hơn mức độ thân thiện giữa các loài,
các nhà sinh vật học muốn lượng hóa quan hệ này.

Họ nhận thấy mức độ thân thiện giữa hai loài liên quan đến \(k\) thuộc tính,
tạm đánh số là thuộc tính \(1,2,\ldots,k\). Mỗi thuộc tính đều có thể đo bằng
một giá trị số. Kết quả nghiên cứu cho thấy: với \(k-1\) thuộc tính đầu, hai
loài càng khác nhau thì càng thân thiện; riêng thuộc tính thứ \(k\) thì ngược
lại, hai loài càng giống nhau thì càng thân thiện.

Vì vậy họ dự đoán có thể dùng công thức sau để đo mức độ thân thiện giữa hai
loài \(a,b\):
\[
  \operatorname{Friendliness}(a,b)=
  \sum_{i=1}^{k-1} C_i\lvert A_{a,i}-A_{b,i}\rvert
  - C_k\lvert A_{a,k}-A_{b,k}\rvert,
\]
trong đó \(C_i\) là các hằng số không âm, còn \(A_{a,i}\) là giá trị thuộc
tính thứ \(i\) của loài \(a\). Khi biết đầy đủ các thuộc tính của từng loài,
ta có thể tính trực tiếp mức độ thân thiện của một cặp loài. Câu hỏi là: trong
các loài đã phát hiện, cặp loài thân thiện nhất là cặp nào, và mức độ thân
thiện của chúng bằng bao nhiêu?

\section{Mô hình toán học}

Đặt \(A_{j,i}\) là giá trị thuộc tính thứ \(i\) của loài \(j\). Nhiệm vụ là
tìm hai chỉ số
\[
  1\le a<b\le n
\]
sao cho hàm mục tiêu
\begin{equation}
  \operatorname{Friendliness}(a,b)=
  \sum_{i=1}^{k-1} C_i\lvert A_{a,i}-A_{b,i}\rvert
  - C_k\lvert A_{a,k}-A_{b,k}\rvert
\label{eq:friendliness}
\end{equation}
đạt giá trị lớn nhất.

\section{Mâu thuẫn của thuật toán trực tiếp}

Thuật toán nguyên thủy nhất là liệt kê mọi cặp \(a,b\), rồi tính giá trị lớn
nhất của \(\operatorname{Friendliness}(a,b)\). Độ phức tạp thời gian khi đó là
\(O(n^2k)\), rõ ràng không thể vượt qua toàn bộ dữ liệu.

Một điều kiện rất dễ chú ý là \(k\le 5\). Điều này gợi ý rằng ta nên nghĩ đến
các thuật toán có thừa số mũ theo \(k\), chẳng hạn \(O(n2^k)\), \(O(nk!)\) hoặc
\(O(n\log n)\) sau khi nhân với một hằng số phụ thuộc vào \(k\).

\section{Thử tối ưu bằng cấu trúc dữ liệu}

Trước hết bỏ dấu trị tuyệt đối trong \eqref{eq:friendliness}. Với mỗi chiều, dấu phía trước hiệu
của hai thuộc tính có thể là dương hoặc âm. Ta thu được dạng
\begin{equation}
  \sum_{i=1}^{k} s_i C_i(A_{b,i}-A_{a,i}),
  \qquad s_i\in\{-1,1\},
\label{eq:signed}
\end{equation}
trong đó dấu được chọn sao cho phù hợp với quan hệ lớn nhỏ của hai giá trị
thuộc tính tương ứng. Do đó có \(2^k\) trường hợp cần xét.

Với mỗi loài \(a\), dùng một số nhị phân \(BN\in[0,2^k-1]\) để biểu diễn cách
chọn dấu trong (2). Nếu bit thứ \(i\) của \(BN\) bằng \(1\), đặt
\(\operatorname{SIGN}(BN,i)=1\); ngược lại đặt
\(\operatorname{SIGN}(BN,i)=-1\). Định nghĩa
\[
  F(a,BN)=\sum_{i=1}^{k}\operatorname{SIGN}(BN,i)C_iA_{a,i}.
\]

Khi đó giá trị lớn nhất có thể viết dưới dạng
\begin{equation}
  \max\{F(b,BN)-F(a,BN)\},
\label{eq:maxf}
\end{equation}
trong đó \(a\ne b\), và với mọi \(i\in[1,k]\), dấu trong \(BN\) phải tương
thích với thứ tự giữa \(A_{a,i}\) và \(A_{b,i}\):
\begin{equation}
  \operatorname{SIGN}(BN,i)(A_{b,i}-A_{a,i})\ge 0.
\label{eq:strict}
\end{equation}

Nếu cài đặt đúng điều kiện \eqref{eq:strict}, với mỗi \(a\) và mỗi \(BN\) ta phải tìm một
loài \(b\) thỏa điều kiện này và có giá trị \(F(b,BN)\) lớn nhất. Có thể nghĩ
đến việc dùng \(2^k\) cây dữ liệu nhiều tầng: tầng theo các thuộc tính, phần
tử lưu trữ là giá trị lớn nhất của \(F\). Khi đó mỗi lần tìm kiếm và chèn có
độ phức tạp \(O(\log^k n)\), tổng thời gian khoảng
\[
  O(n2^k\log^k n),
\]
còn bộ nhớ khoảng \(O(n\log^{k-1}n)\). Dù đã tốt hơn liệt kê cặp, cách này vẫn
có nhiều mâu thuẫn: thời gian còn cao, bộ nhớ lớn, và độ phức tạp cài đặt rất
khó chấp nhận.

\section{Nới điều kiện}

Nguyên nhân căn bản của các mâu thuẫn trên là yêu cầu \eqref{eq:strict} quá khắt khe. Ta thử
nới yêu cầu này.

Nếu bỏ hẳn \eqref{eq:strict}, kết quả sẽ sai, vì khi đó việc chọn dấu là tùy tiện. Tuy nhiên
thuộc tính thứ \(k\) có tính chất đặc biệt: nó đi vào hàm mục tiêu với dấu trừ,
còn \(k-1\) thuộc tính đầu đi vào với dấu cộng. Vì vậy chỉ cần giữ đúng thứ tự
ở thuộc tính thứ \(k\); các thuộc tính còn lại có thể để việc liệt kê \(BN\)
tự chọn dấu tốt nhất.

Cụ thể, thay \eqref{eq:strict} bằng điều kiện yếu hơn:
\begin{equation}
  \operatorname{SIGN}(BN,k)(A_{b,k}-A_{a,k})\ge 0.
\label{eq:relaxed}
\end{equation}

Thực nghiệm với thuật toán nguyên thủy cho kết quả hoàn toàn giống nhau. Phần
chứng minh dưới đây cho thấy cách nới điều kiện này là đúng.

\paragraph{Định lý.}
Với mọi \(A,B\in\mathbb{R}\),
\[
  \lvert A-B\rvert\ge A-B,\qquad
  \lvert A-B\rvert\ge B-A.
\]

\paragraph{Hệ quả.}
Với mọi cặp loài \(a,b\), nếu cố định dấu của chiều thứ \(k\) đúng với thứ tự
giữa \(A_{a,k}\) và \(A_{b,k}\), thì việc lấy cực đại trên mọi lựa chọn dấu ở
các chiều \(1,\ldots,k-1\) cho giá trị đúng bằng
\[
  \sum_{i=1}^{k-1} C_i\lvert A_{a,i}-A_{b,i}\rvert
  - C_k\lvert A_{a,k}-A_{b,k}\rvert.
\]

Gọi \(S_1\) là tập các bộ \((a,b,BN)\) thỏa điều kiện đầy đủ \eqref{eq:strict}, còn \(S_2\)
là tập các bộ thỏa điều kiện nới lỏng \eqref{eq:relaxed}. Rõ ràng \(S_1\subseteq S_2\), nên
nếu
\[
  M_1=\max_{(a,b,BN)\in S_1} \{F(b,BN)-F(a,BN)\},\qquad
  M_2=\max_{(a,b,BN)\in S_2} \{F(b,BN)-F(a,BN)\},
\]
thì \(M_1\le M_2\).

Mặt khác, với bất kỳ \((a,b,BN_2)\in S_2\), nhờ định lý trên, ta luôn có thể
chọn lại các dấu ở \(k-1\) chiều đầu để thu được một \(BN_1\) thỏa \eqref{eq:strict}, đồng
thời giá trị \(F(b,BN_1)-F(a,BN_1)\) không nhỏ hơn giá trị tương ứng của
\(BN_2\). Do đó \(M_1\ge M_2\). Suy ra \(M_1=M_2\).

Vậy chỉ cần bảo đảm điều kiện \eqref{eq:relaxed}, cực đại của \eqref{eq:maxf} chính là cực đại của
\(\operatorname{Friendliness}\).

\section{Cài đặt}

Sắp xếp các loài theo khóa \(A_{\cdot,k}\) tăng dần. Sau khi sắp xếp, chỉ cần
xét các cặp trong đó loài đứng trước có giá trị thuộc tính thứ \(k\) không lớn
hơn loài đứng sau.

Với mỗi mẫu dấu \(BN\) có bit thứ \(k\) bằng \(1\), duy trì
\[
  \operatorname{Best}(BN)=\max F(b,BN)
\]
trên các loài \(b\) đã quét qua. Khi xét loài hiện tại \(a\), dùng
\[
  \operatorname{Best}(BN)-F(a,BN)
\]
để cập nhật đáp án, rồi đưa \(F(a,BN)\) vào \(\operatorname{Best}(BN)\).
Làm tương tự cho các mẫu dấu đã bao phủ mọi lựa chọn ở \(k-1\) thuộc tính đầu.

Việc tính tất cả các giá trị \(F(a,BN)\) cần \(O(n2^k)\) thời gian. Phần sắp
xếp cần \(O(n\log n)\). Do đó độ phức tạp tổng là
\[
  O(n\log n+n2^k),
\]
và bộ nhớ là \(O(n2^k)\), hoặc \(O(2^k)\) nếu tính \(F\) theo dòng trong quá
trình quét.

\section{Kết luận}

Các bước giải bài toán bằng phương pháp ``nới điều kiện'' có thể khái quát như
sau:

\begin{enumerate}
  \item Chuyển nhiệm vụ thành mô hình toán học.
  \item Xây dựng thuật toán nguyên thủy.
  \item Thử tối ưu thuật toán nguyên thủy bằng cấu trúc dữ liệu.
  \item Phân tích điều kiện khắt khe sau khi tối ưu, rồi chuyển nó thành điều
  kiện yếu hơn nhưng vẫn giữ nguyên giá trị tối ưu.
\end{enumerate}

\end{document}
